\subsection{分解与Grothendieck群}

若 \(\mathcal{C}\) 是由A-模的类组成的一个集合，我们将称左分解 \((\mathbf{P}, p)\) 是 \(\mathcal{C}\) 型有界的，如果复形 \(\mathbf{P}\) 是 \(\mathcal{C}\) 型有界的 (X，第41页)。

定理1. --- 设 \(\mathcal{C}_{0}\) 与 \(\mathcal{C}\) 是两个加性且左正合的A-模类，使得 \(\mathcal{C}_{0} \subset \mathcal{C}\)，并且每个 \(\mathcal{C}\) 型的A-模都具有 \(\mathcal{C}_{0}\) 型的有界左分解。则同态 \(\alpha : \mathrm{K}(\mathcal{C}_{0}) \to \mathrm{K}(\mathcal{C})\)（由 \(\mathcal{C}_{0}\) 到 \(\mathcal{C}\) 的包含关系导出的）是双射；若 M 是一个 \(\mathcal{C}\) 型的A-模，且 P 是 M 的一个 \(\mathcal{C}_{0}\) 型有界左分解，则我们有 \(\alpha^{-1}([M]_{\mathcal{C}}) = \chi_{\mathcal{C}_{0}}(P)\) （X，第 41 页，例 6）。

引理4. --- 设 \(f: \mathbf{M}' \to \mathbf{M}\) 是 \(\mathcal{C}\) 型A-模的同态，且 \(p: \mathbf{P} \to \mathbf{M}\) 是 \(\mathbf{P}\) 的一个 \(\mathcal{C}_0\) 型有界的左分解。存在一个左分解 \(p': \mathbf{P}' \to \mathbf{M}'\)（\(\mathcal{C}_0\) 型有界的）以及一个复形的态射 \(u: \mathbf{P}' \to \mathbf{P}\) ，使得 \(p \circ u = f \circ p'\)。

让我们对P的长度n进行归纳推理，当该长度为 \textless{} 0时命题是平凡的。考虑映射 \(g: \mathbf{M}' \times \mathbf{P}_0 \to \mathbf{M}\) 使得

\[
g (x, y) = f (x) - p _ {0} (y) \quad \text { for } \quad x \in \mathbf {M} ^ {\prime}, y \in \mathbf {P} _ {0},
\]

及其核K；由于g是满射且 \(M' \times P_{0}\) 和M是C类型的，A-模K是C类型的。设 \(h : P_{0}' \to K\) 是一个满射同态，其中 \(P_{0}'\) 是 \(C_{0}\) 型的；记 \(p_{0}' : P_{0}' \to M'\)（相应地， \(u_{0} : P_{0}' \to P_{0}\)）为 h 与投影 \(K \to M\)（相应地， \(K \to P_{0}\)）的复合同态；同态 \(p_{0}'\) 是满射的，并且我们有一个交换图

\[
\begin{array}{c} \mathbf {P} _ {0} ^ {\prime} \xrightarrow {u _ {0}} \mathbf {P} _ {0} \\ p _ {0} ^ {\prime} \Big \downarrow \qquad \qquad \qquad \qquad \qquad \qquad \qquad \qquad \qquad \Big \downarrow p _ {0} \\ \mathbf {M} ^ {\prime} \xrightarrow [ f ]{} \mathbf {M}. \end{array}
\]

然后只需将归纳假设应用于同态

\[
\operatorname{Ker} p _ {0} ^ {\prime} \rightarrow \operatorname{Ker} p _ {0}
\]

由 \(u_{0}\) .

引理5. --- 考虑一个交换图

\[
\begin{array}{c} \mathbf {P} ^ {\prime} \xrightarrow {u} \mathbf {P} \\ p ^ {\prime} \Big \downarrow \quad \Big \downarrow p \\ 0 \longrightarrow \mathbf {M} ^ {\prime} \xrightarrow {f} \mathbf {M} \longrightarrow \mathbf {M} ^ {\prime \prime} \longrightarrow 0 \end{array}
\]

其中 (P, p)（相应地 (P′, p′)）是 M（相应地 M′）的左分解，且底部水平行是一个正合序列。存在一个拟同构 \(p''\) : Con \((u) \to \mathbf{M}''\) 。

事实上，正合序列 (X，第 37 页，命题 7)

\[
0 \to \mathrm{P} \xrightarrow {\pi} \operatorname{Con} (u) \xrightarrow {\delta} \mathrm{P} ^ {\prime} (- 1) \to 0
\]

给出一个正合同调序列

\[
\begin{array}{c}\rightarrow \mathrm{H} _ {n} (\mathrm{P}) \rightarrow \mathrm{H} _ {n} (\mathrm{Con} (u)) \rightarrow \mathrm{H} _ {n - 1} (\mathrm{P} ^ {\prime}) \rightarrow \dots\\\dots \rightarrow \mathrm{H} _ {1} (\mathrm{Con} (u)) \rightarrow \mathrm{H} _ {0} (\mathrm{P} ^ {\prime}) \stackrel {{\partial}} {{\rightarrow}} \mathrm{H} _ {0} (\mathrm{P}) \rightarrow \mathrm{H} _ {0} (\mathrm{Con} (u)) \rightarrow 0.\end{array}
\]

根据 X，第 38 页，引理 3 a)，我们有 \(\partial = -\mathrm{H}_{0}(u)\) 。由于 \(\mathrm{H}_{n}(\mathbf{P}) = 0 = \mathrm{H}_{n}(\mathbf{P}')\)（\(n > 0\)），且 \(\mathrm{H}_{0}(u): \mathrm{H}_{0}(\mathrm{P}') \to \mathrm{H}_{0}(\mathbf{P})\) 与 \(f: \mathbf{M}' \to \mathbf{M}\) 等同，我们得出结论 \(\mathrm{H}_{n}(\mathrm{Con}(u)) = 0\)（对 \(n > 0\)），并且 \(\mathrm{H}_{0}(\mathrm{Con}(u))\) 同构于 \(\mathbf{M}''\) ，从而引理得证。

现在让我们证明该定理。

\begin{enumerate}
\def\labelenumi{\alph{enumi})}
\tightlist
\item
  设 M 是一个 C 型的 A-模。对于 M 的每一个左分解 \((\mathbf{P}, p)\)（\(C_{0}\) 型有界的），\(\chi_{\mathcal{C}_{0}}(\mathbf{P})\) 作为 \(\mathbf{K}(\mathcal{C}_{0})\) 中的元素仅依赖于 M。事实上，设 \((\mathbf{P}_{1}, p_{1})\) 与 \((\mathbf{P}_{2}, p_{2})\) 是该类型的两个分解。考虑 A-模
\end{enumerate}

\[
(\mathbf {P} _ {1} \times \mathbf {P} _ {2}, p _ {1} \times p _ {2})
\]

是A-模 \(\mathbf{M} \times \mathbf{M}\) 的分解，以及同态 \(\Delta : x \mapsto (x, x)\)（从 \(\mathbf{M}\) 到 \(\mathbf{M} \times \mathbf{M}\)）。根据引理 4，存在一个分解 \((\mathbf{Q}, q)\)（\(\mathbf{M}\) 的 \(\mathcal{C}_0\) 型有界分解）和一个交换图

\[
\begin{array}{c c c} \mathbf {Q} & \stackrel {{u}} {{\longrightarrow}} \mathbf {P} _ {1} & \times \mathbf {P} _ {2} \\ q \Big \downarrow & & \Big \downarrow p _ {1} \times p _ {2} \\ \mathbf {M} & \stackrel {{\Delta .}} {{\longrightarrow}} \mathbf {M} & \times \mathbf {M}; \end{array}
\]

我们推出一个交换图

\[
\begin{array}{c c c} \mathrm{Q} & \xrightarrow {u \circ p r _ {i}} & \mathrm{P} _ {i} \\ q \Big \downarrow & & \Big \downarrow p _ {i} \\ \mathbf {M} & \xrightarrow {\mathrm{l} _ {\mathrm{M}}} & \mathbf {M}, \end{array} \qquad i = 1, 2.
\]

根据引理 5，Con \((u \circ pr_i)\) 具有零同调，因此 \(u \circ pr_i\) 是一个拟同构且 \(\chi_{\mathcal{C}_0}(\mathbf{Q}) = \chi_{\mathcal{C}_0}(\mathbf{P}_i)\) (X，第 41 页，命题 10)；由此得出 \(\chi_{\mathcal{C}_0}(\mathbf{P}_1) = \chi_{\mathcal{C}_0}(\mathbf{P}_2)\) 如所宣称。

\begin{enumerate}
\def\labelenumi{\alph{enumi})}
\setcounter{enumi}{1}
\tightlist
\item
  对于 C 型的每个A-模 M，设 \(\varphi(M) \in K(\mathcal{C}_{0})\) 为 \(\chi_{\mathcal{C}_{0}}(P)\) 对 M 的所有 \(C_{0}\) 型有界左分解 P 所取的公共值。让我们证明函数 \(\varphi : \mathcal{C} \to K(\mathcal{C}_{0})\) 是可加的。为此设
\end{enumerate}

\[
0 \to \mathbf {M} ^ {\prime} \xrightarrow {f} \mathbf {M} \to \mathbf {M} ^ {\prime \prime} \to 0
\]

设为一个C型A-模的正合序列。根据引理4，存在一个交换图

\[
\begin{array}{c} \mathrm{P} ^ {\prime} \xrightarrow {u} \mathrm{P} \\ p ^ {\prime} \Big \downarrow \qquad \qquad \qquad \Big \downarrow p \\ 0 \longrightarrow \mathrm{M} ^ {\prime} \xrightarrow {f} \mathrm{M} \longrightarrow \mathrm{M} ^ {\prime \prime} \longrightarrow 0 \end{array}
\]

其中 (P, p) 和 (P', p') 是 \(\mathcal{C}_0\) 型的有界左分解。于是我们有

\[
\varphi (\mathbf {M}) = \chi_ {\mathcal {C} _ {0}} (\mathbf {P}), \quad \varphi (\mathbf {M} ^ {\prime}) = \chi_ {\mathcal {C} _ {0}} (\mathbf {P} ^ {\prime})
\]

并且由引理5

\[
\varphi (\mathbf {M} ^ {\prime \prime}) = \chi_ {\mathcal {C} _ {0}} (\operatorname{Con} (u)) = \chi_ {\mathcal {C} _ {0}} (\mathrm{P}) - \chi_ {\mathcal {C} _ {0}} (\mathrm{P} ^ {\prime}) = \varphi (\mathrm{M}) - \varphi (\mathrm{M} ^ {\prime});
\]

这正是我们想要证明的。

\begin{enumerate}
\def\labelenumi{\alph{enumi})}
\setcounter{enumi}{2}
\tightlist
\item
  设此时 \(\beta: \mathrm{K}(\mathcal{C}) \to \mathrm{K}(\mathcal{C}_0)\) 是这样一个同态，使得在上述记号下，我们有 \(\beta([\mathrm{M}]_{\mathcal{C}}) = \chi_{\mathcal{C}_0}(\mathrm{P})\) 。由于 \(p\) 是一个拟同构，我们有 \(\chi_{\mathcal{C}}(\mathrm{P}) = [\mathrm{M}]_{\mathcal{C}}\) ，因此 \(\alpha \circ \beta([\mathrm{M}]_{\mathcal{C}}) = \alpha(\chi_{\mathcal{C}_0}(\mathrm{P})) = \chi_{\mathcal{C}}(\mathrm{P}) = [\mathrm{M}]_{\mathcal{C}}\) ，故 \(\alpha \circ \beta = 1_{\mathrm{K}(\mathcal{C})}\) 。如果 M 是 \(\mathcal{C}_0\) 型的，那么 (M, \(1_{\mathrm{M}}\) ) 是 M 的一个分解，因此 \(\varphi(\mathrm{M}) = [\mathrm{M}]_{\mathcal{C}_0}\) ，且 \(\beta \circ \alpha = 1_{\mathrm{K}(\mathcal{C}_0)}\) 也成立，这就完成了证明。
\end{enumerate}

我们将在 § 8（X，第 137 页）中将此定理应用于“有限投射维数”的模。

